\section{Results}
\label{sec:results}

\subsection{Local refinement retains q4 work at a physical crossing}
\label{sec:local}
The 12-block A5000 confirmation directly tests local lowering at two fixed 32B histories (Figure~\ref{fig:confirmation}). The noncrossing primary endpoint has a 0.43\% point reduction in verification-through-commit time, with a 95\% interval of $[-1.13\%,2.01\%]$. At the selected crossing supporting endpoint, LOCAL reduces time by 21.56\% [20.21\%, 22.83\%]. It retains q4 work otherwise split into q3+q1; both policies keep width four in the noncrossing history. The supporting crossing/noncrossing interaction is 0.788 [0.768, 0.806], a ratio of paired time ratios with null value one. These intervals describe repeated time blocks for the fixed histories.

The separate 8B target-only study in the right panel of Figure~\ref{fig:confirmation} has a primary equal-conversation reduction of 1.44\% $[-0.82\%,3.85\%]$ under its response-plus-reset clock. Section~\ref{sec:natural} examines the accompanying grouping and consumer costs. Its two fixed conversations and the 32B histories are analyzed separately; neither primary interval resolves a benefit direction.

\paperfigure[!htbp]{f07_confirmation}{LOCAL/SIMPLE effects in 12 paired A5000 time blocks. Left: lines connect the same block's two fixed 32B histories (one root; C4 on, P4 off, frozen proposals; verification through commit). Noncrossing is primary; crossing and their ratio-of-ratios interaction are supporting. Right: the separately analyzed 8B target-only primary endpoint weights two fixed conversations equally in log time, from response through reset. Positive $100(1-\mathrm{LOCAL}/\mathrm{SIMPLE})\%$ means less LOCAL time. Printed limits give individual 95\% joint whole-block intervals. The interaction has null value one.}{Two panels show all twelve numbered time blocks. Lines connect cases within each block. The 32B crossing effects consistently favor LOCAL, while the primary noncrossing interval spans zero. The separately clocked 8B target-only panel retains both conversations and their equal-weight diamond; its primary interval also spans zero. Printed intervals accompany the estimates, and the supporting ratio-of-ratios interaction is below one.}{fig:confirmation}

The original two-repeat study supplies an independently clocked explanation of this retained work. At the crossing, projection event time falls from 74.334 to 53.417\,ms as q3+q1 becomes q4; attention event time changes from 12.060 to 11.621\,ms, with 0.134\,ms of added concatenations. These device events locate work inside the lowering and have a separate clock from transaction wall time. The treatment changes widths, fusion, launches and views together. Section~\ref{supp:local} reports the original paired times and diagnostics.

Each measured first transaction continues from its own committed state, exercising the owned-row, trim and successor-view connection in Section~\ref{sec:method}. The direct evidence thus joins retained projection work to a valid complete transaction. Independent request studies next measure few-column execution and prefill packing under their own controls.

\subsection{Few-column arithmetic has value beyond scalar regrouping}
\label{sec:c4}
The RTX confirmation separates scalar dispatch organization from the admitted few-column path using the B/C1/C4 controls in Table~\ref{tab:dispatch}, with SIMPLE fixed and P4 off. Its sole primary comparison, C1/B, has a 0.90\% point increase in resident-request time, with a 95\% interval from a 0.48\% reduction to a 2.19\% increase. The clock runs from reset through final synchronization. Mean target verification saves 100.43\,ms, while target prefill adds 173.57\,ms and drafting adds 7.43\,ms. The unrounded phase values, including remaining stages, give an 80.52\,ms mean service increase. Figure~\ref{fig:c4} separates this arithmetic account from the geometric request effect.

The descriptive comparisons favor C4: request time falls by 15.80\% relative to C1 [13.07\%, 18.30\%] and by 15.04\% relative to B [12.40\%, 17.49\%]. All 24 root aggregates favor C4 against either scalar control. Under this target profile, few-column execution reduces complete-request time beyond the change obtained from scalar regrouping.

\paperfigure[!htbp]{f04_few_column}{RTX resident-request controls with SIMPLE fixed and P4 off (reset through final synchronization). Dots show 24 root effects per contrast; diamonds and bars give equal-root/equal-block geometric reductions and individual 95\% crossed stratified-root/whole-block intervals. C1/B is the sole primary contrast; C4 contrasts are descriptive. Positive values mean time saved. The waterfall instead reports arithmetic mean milliseconds saved, B minus C1; it does not decompose the geometric effects.}{A dispatch matrix distinguishes scalar loop B, grouped scalar C1 and few-column C4. Each contrast retains twelve code and twelve math root effects plus its geometric summary. The primary C1/B interval crosses zero; both descriptive C4 intervals favor the candidate. A separate arithmetic waterfall shows verification savings outweighed by prefill cost, ending at a net cost of 80.52 milliseconds. All signed values and the zero lines remain visible.}{fig:c4}

C4 records 2,273,400 verification admissions and 2,652 prefill admissions, the latter from one root's q3 tail in each of its 12 requests. These exposures define the measured target few-column treatment. Section~\ref{supp:rtx-confirmation} gives the corresponding root, block and phase estimates.

The independent A5000 C1/C4 study records a 27.77\% request reduction [25.43\%, 30.09\%] across 12 roots and two rounds, with all few-column admissions in verification. Figure~\ref{fig:c4-initial} and Table~\ref{tab:direct-effects} retain that study and the earlier two-round RTX results under their own controls and estimators. Together they quantify the request value of admitted groups; their execution remains SIMPLE with P4 off.

\FloatBarrier
\subsection{Prefill organization makes several admitted groups useful}
\label{sec:p4}
Independent P4 studies test organization between complete groups in wide prefill, with SIMPLE fixed. Their four arms distinguish the packing increment within target-only or speculative execution from speculation's value with packing already enabled (Figure~\ref{fig:p4}). The two deployments contribute 12 and 24 study-specific roots; each root appears in all three contrasts, giving 108 effect marks. Prompt-length order supplies a readable root rank, not a length-scaled or causal axis.

\paperfigure[!htbp]{f05_prefill}{Independent A5000 and RTX P4 controls with SIMPLE fixed (resident requests, reset through final synchronization). The studies contain 12 and 24 roots, each reused in three contrasts (36 study-specific roots; 108 effect marks). Horizontal positions follow prompt-length rank, not proportional token spacing; the studies are unpaired. Row labels specify candidate/control; positive reductions mean less candidate time. Diamonds and bars give equal-root geometric reductions and individual 95\% code/math-stratified root intervals. The outlined A5000 \texttt{train/279} effect is $-0.33\%$.}{Three contrast rows appear in two independent deployment columns. Every root effect is retained at a rank ordered by prompt length, with a separate overall diamond and confidence interval. Packing favors the candidate for every root in the first two rows. At fixed P4, the A5000 train/279 point is an outlined diamond at a small negative reduction; it is not moved to zero. The two studies are not paired.}{fig:p4}

A5000 P4 reduces target-only request time by 13.18\% [10.38\%, 15.73\%] and speculative request time by 14.95\% [11.80\%, 17.76\%] relative to C4. The corresponding independent RTX reductions are 46.19\% [42.72\%, 49.32\%] and 49.53\% [47.00\%, 51.87\%]. Each contrast uses the contemporaneous unpacked control in the same decoding policy.

With P4 fixed, speculation reduces time relative to target-only execution by 13.60\% on A5000 [11.26\%, 15.89\%] and 14.91\% on RTX [12.50\%, 17.44\%]. A5000 \texttt{train/279} is 0.33\% slower; the other eleven roots are faster. This same-study comparison determines whether live speculation repays its additional work after prefill has been improved.

The A5000 packing gain is concentrated in prefill organization: it saves 1,064.5\,ms of the 1,077.1\,ms net arithmetic mean speculative-request saving. All 232 P4-entering prefill calls use aligned q32 groups within one physical extent. RTX speculative prefill falls from 6.852 to 2.395\,s. The complete phase accounts and observed dispatch widths are retained in Section~\ref{supp:effects}.

\FloatBarrier
\subsection{Natural opportunities and service costs}
\label{sec:natural}
A selected crossing exposes work that a natural request may rarely need. M1 measures this frequency in its WHOLE/LOCAL $\times$ loop4/packed factorial (Table~\ref{tab:m1}). Each arm has 2,224 logical cycles. WHOLE executes 2,244 physical verification calls and LOCAL 2,224, removing 20 extra calls, or 0.90\% of logical cycles.

\papertable[!htbp]{t05_m1}{M1 factorial: 16 roots and four sessions per root, with 64 measured requests per cell. Each projection row reports native-service totals and its direct LOCAL/WHOLE time reduction. Intervals are 99\% marginal members of the five-comparison Bonferroni family; Table~\ref{tab:m1-family} retains the complete family.}{tab:m1}

At fixed loop4, LOCAL reduces total native-service time by 0.209\%, with 99\% interval [0.042\%, 0.463\%]. At fixed packed projection, the reduction is 0.128\%, with interval $[-0.057\%,0.405\%]$. Packing improves both cut policies, while the root-geometric interaction is 1.00070 [1.00007, 1.00139], indicating mild attenuation of LOCAL's effect under packing. The small service effects occur with the measured sparse crossing geometry.

\begin{samepage}
The independent 8B target-only confirmation measures the original two conversations over 12 time blocks. Its primary equal-conversation estimate is a 1.44\% time reduction, with 95\% interval $[-0.82\%,3.85\%]$ (Table~\ref{tab:confirmation-natural}). Conversation 0 retains identical physical grouping; Conversation 1 removes one outer call, 273 to 272, at equal semantic work. Their mean service differences are $-95.8$ and $-190.0$\,ms, while common sampling differences are $-112.6$ and $-192.9$\,ms. Common consumer variation is comparable to the net effect even without a grouping change.\par
\end{samepage}

\papertable[htbp]{t07_confirmation_natural}{Natural 8B target-only confirmation: effects and individual 95\% whole-time-block intervals over two fixed conversations and 12 paired blocks. The primary estimate weights conversations equally within each block. All block values and direction counts remain in Section~\ref{supp:confirmation}.}{tab:confirmation-natural}

The earlier natural panels show higher time under LOCAL: 0.73\% for historical RTX live speculation and 4.22\% for the two-round A5000 target-only study (Table~\ref{tab:natural}). In the latter, common sampling contributes 1.420\,s of the 1.535\,s increase. Its ratio-of-totals estimator and the RTX panel's original estimator remain separate from the confirmation. Interpreting service changes therefore requires both the execution geometry and the accompanying request costs.

\FloatBarrier
\subsection{Draft and setup costs can outweigh target savings}
\label{sec:cost}
In the formal B570 study, compatible grouping reduces target-only time by 21.73\% [20.35\%, 23.01\%] and speculative time by 12.78\% [12.56\%, 12.99\%], across 24 roots and two rounds. Complete speculative service nevertheless takes 2.512 times its grouped target-only control [2.486, 2.539], with all 24 roots slower (Figure~\ref{fig:b570}). Grouped target work is useful in this CPU-draft deployment, but target-only generation has the lower service cost.

\paperfigure[!htbp]{f06_b570}{Grouping and complete speculative service in the formal B570 CPU-draft study (24 roots, two rounds). Each comparison reuses the same roots (72 effect marks); open diamonds give equal-root estimates with individual 95\% stratified-root intervals printed at right. Ratios share a logarithmic axis: one means equal time, and lower values mean less candidate time. Service includes CPU drafting and normal teardown.}{Three rows each retain 24 measured root ratios and an open-diamond equal-root summary on one continuous log axis. Grouped/scalar target-only execution is 0.783 with 95\% interval 0.770 to 0.797; grouped/scalar speculation is 0.872 with interval 0.870 to 0.874. Grouped speculation versus grouped target-only is 2.512 with interval 2.486 to 2.539, to the right of the equal-time line at one and therefore slower. The 72 root marks reuse the same 24 roots across three contrasts. Numeric limits retain the narrow grouping intervals without changing their graphical extent.}{fig:b570}

\FloatBarrier
The separate two-root draft-placement pilot takes 2.326 times the grouped target-only time after moving the main draft matrices and LM head to Vulkan0. Its drafting stages average 14.807\,s against 1.643\,s saved in target work, meeting the predeclared service stop (Table~\ref{tab:pilot}). A historical Llama study likewise reverses a resident-request advantage when loading, checks and primers enter the short-worker lifetime (Table~\ref{tab:historical-applications}). Draft placement and setup lifetime can therefore change the preferred decoding policy even when the target kernels become faster.

